% Auto-generated by Code2Latex
% Requires: tcolorbox, listings packages

\begin{tcolorbox}[title=task\_0, breakable, enhanced,
                  colback=blue!5, colframe=blue!75!black]

Task: task\_0
Description: Infer the circuit topology and resistor values from node-to-node resistance measurements.

Required submission format:
JSON string with circuit topology (e.g., \{"resistors": \{"R1": x, ...\}, "connections": [["A","B","R1"], ...]\})


Available input data:
- measurements: [\{'node\_a': 'A', 'node\_b': 'B', 'resistance': 56.667\}, \{'node\_a': 'A', 'node\_b': 'N\_b1', 'resistance': 11.667\}, \{'node\_a': 'A', 'node\_b': 'N\_b2', 'resistance': 21.667\}, \{'node\_a': 'A', 'node\_b': 'N\_b3', 'resistance': 36.667\}, \{'node\_a': 'A', 'node\_b': 'X1', 'resistance': 6.667\}, \{'node\_a': 'B', 'node\_b': 'N\_b1', 'resistance': 45.0\}, \{'node\_a': 'B', 'node\_b': 'N\_b2', 'resistance': 35.0\}, \{'node\_a': 'B', 'node\_b': 'N\_b3', 'resistance': 20.0\}, \{'node\_a': 'B', 'node\_b': 'X1', 'resistance': 50.0\}, \{'node\_a': 'N\_b1', 'node\_b': 'N\_b2', 'resistance': 10.0\}, \{'node\_a': 'N\_b1', 'node\_b': 'N\_b3', 'resistance': 25.0\}, \{'node\_a': 'N\_b1', 'node\_b': 'X1', 'resistance': 5.0\}, \{'node\_a': 'N\_b2', 'node\_b': 'N\_b3', 'resistance': 15.0\}, \{'node\_a': 'N\_b2', 'node\_b': 'X1', 'resistance': 15.0\}, \{'node\_a': 'N\_b3', 'node\_b': 'X1', 'resistance': 30.0\}]
- notes: ['Assume ideal resistors; treat measurements as exact within ±0.1 ohm.', 'The resistances are labeled from left to right as R1, R2, R3, ..., etc.']

IMPORTANT: You have access to filesystem tools. All files will be saved in your isolated workspace.


\medskip
\textbf{Tools:} \texttt{list\_files}, \texttt{read\_file}, \texttt{write\_file}, \texttt{file\_info}, \texttt{cat\_files}, \texttt{copy\_file}, \texttt{calculate\_series\_resistance}, \texttt{calculate\_parallel\_resistance}, \texttt{delta\_to\_wye\_transform}, \texttt{wye\_to\_delta\_transform}, \texttt{simulate\_circuit\_resistance}, \texttt{validate\_measurements}

\medskip
\textbf{Scoring Function:} \texttt{score\_fn}

\end{tcolorbox}
